\scanpage{49}
\uline{Résumé} On peut \unsure{exprimer|définir} des V.A. sur $k$ (\unsure{alg|alg.} clos fixé), \deleted{\unsure{des}} en termes de cycles algébriques \inserted{alg. équiv. à $0$} sur une variété, soit via les diviseurs ($\deleted{\unsure{A}}B_X$\inserted{$={}$ Picard \unsure{connexe|réduit}}) soit via les $0$-cycles ($A_X\deleted{\unsure{)}}={}$ Albanese). On obtient \deleted{donc}\inserted{alors}, pour deux variétés abéliennes $A$ et $B$, définies \inserted{chacune,} en termes de $X$ \deleted{resp.} de $Y$, de l'une ou de l'autre façon, et un cycle alg. de dim convenable $Z$ sur $X\times Y$, par $\unsure{\mathscr{D}|\delta|D}\unsure{\leadsto|-|\mapsto} Z(\unsure{\delta|\mathscr{D}|D})$, une application $A(k)\longrightarrow B(k)$, \deleted{\unsure{et une|et même}} (il y a env. $4=2\times2$ cas à distinguer). Nous montrons qu'elle est tjrs associée à un hom de V.A. $u_Z:A\longrightarrow B$, \inserted{\unsure{ne|en} dépendant que de la classe d'\unsure{équiv|équivalence} \uline{numérique} de $Z$.} \deleted{et \unsure{pour|que}}\inserted{D'ailleurs,} si on admet des cycles à coeff. rationnels, \deleted{\unsure{les|tous}} tout hom de V.A. $A\unsure{-|\longrightarrow} B$ s'obtient de cette façon. En termes de cohomologie $l$-adique, on peut dire que les hom $H^1(B)-H^1(A)$ associés à des $\mathrm{Q}$-hom de V.A. $A-B$, sont aussi ceux associés à des $\mathrm{Q}$-cycles
% archive sheet number: 48
